\documentclass[12pt]{amsart}

\usepackage[utf8]{inputenc}
\usepackage[T1]{fontenc}
\usepackage{mathtools}
\usepackage{enumitem}
\usepackage{hyperref}
\usepackage{booktabs}
\usepackage{array}
\usepackage[margin=1in]{geometry}
\usepackage{microtype}

\newtheorem{theorem}{Theorem}[section]
\newtheorem{lemma}[theorem]{Lemma}
\newtheorem{proposition}[theorem]{Proposition}
\newtheorem{corollary}[theorem]{Corollary}
\newtheorem{conjecture}[theorem]{Conjecture}
\theoremstyle{definition}
\newtheorem{definition}[theorem]{Definition}
\newtheorem{example}[theorem]{Example}
\theoremstyle{remark}
\newtheorem{remark}[theorem]{Remark}

\title{Connections Between the Berggren Tree\\
and Millennium Prize Problems}

\author{Paul Klemstine}
\address{Independent Researcher, Menasha, WI}

\date{March 2026}

\begin{document}

\begin{abstract}
Integer factoring and BSD rank computation are Turing-equivalent for
semiprimes: the Selmer group computation for $E_n : y^2 = x^3 - n^2 x$
encodes the factorization of $n = pq$, and conversely, given the factors,
$2$-descent runs in $O(2^{\omega(n)})$ time. We establish this equivalence
and survey further connections between the Berggren Pythagorean triple tree
and the Millennium Prize Problems. For the Riemann Hypothesis, the tree's
prime enrichment ($6.7\times$) and zeta machine provide computational tools
but no proof pathway. For BSD, every PPT maps to a rational point on a
congruent number curve, making the tree an infinite constructor of congruent
numbers. For Hodge, Navier--Stokes, Yang--Mills, and P~vs~NP, we report
negative results from 350+ mathematical fields explored: all approaches to
factoring reduce to five known complexity families, and the Dickman
information barrier is proven fundamental for generic sieves.
\end{abstract}

\maketitle

% ═════════════════════════════════════════════════════════════════════════════
\section{Introduction}\label{sec:intro}

The Berggren Pythagorean triple tree connects to several Millennium Prize
Problems through its algebraic structure, its relationship to elliptic curves,
and its role as a testing ground for complexity-theoretic questions. This
paper presents the connections we have established---both positive results
and honest negative results---across a comprehensive research program
examining 350+ mathematical fields.

% ═════════════════════════════════════════════════════════════════════════════
\section{Factoring and BSD: Turing Equivalence}\label{sec:bsd}

\subsection{Congruent number curve mapping}

\begin{theorem}[Congruent number curve mapping]\label{thm:congruent}
Every PPT $(a,b,c)$ maps to a rational point on
\[
  E_n : y^2 = x^3 - n^2 x, \qquad n = ab/2,
\]
via $x = (c/2)^2$, $y = c(b^2 - a^2)/8$.
\end{theorem}

\begin{proof}
We verify $y^2 = x^3 - n^2 x$. Left side: $c^2(b^2-a^2)^2/64$. Right side:
$c^6/64 - a^2 b^2 c^2/16 = c^2(c^4 - 4a^2 b^2)/64$. Using $a^2 + b^2 = c^2$:
$(b^2-a^2)^2 = c^4 - 4a^2 c^2 + 4a^4$ and
$c^4 - 4a^2 b^2 = c^4 - 4a^2(c^2 - a^2) = c^4 - 4a^2 c^2 + 4a^4$.
These are equal. Verified on 50 triples.
\end{proof}

\begin{corollary}
The Berggren tree generates an infinite family of congruent numbers with
explicit rational points. At depth~$d$, $3^d$ new congruent numbers $n = ab/2$
are produced, each with a constructive witness $(x,y) \in E_n(\mathbb{Q})$.
\end{corollary}

\subsection{The Turing equivalence}

\begin{theorem}[Factoring--BSD rank equivalence]\label{thm:fact-bsd}
For a semiprime $n = pq$, computing $\mathrm{rank}(E_n(\mathbb{Q}))$ by
$2$-descent is Turing-equivalent to factoring~$n$.
\end{theorem}

\begin{proof}[Proof sketch]
For $E_n : y^2 = x(x-n)(x+n)$, $2$-descent tests local solubility of
homogeneous spaces at each prime dividing $2n$. For $n = pq$:
\begin{itemize}
  \item The local condition at~$p$ depends on $n/p = q$.
  \item The local condition at~$q$ depends on $n/q = p$.
\end{itemize}
Thus the Selmer computation IS a factoring computation. Conversely, given
$p$ and $q$, the Selmer bound
$|\mathrm{Sel}^2(E_n)| \le 2^{\omega(2n)+1}$ is computable in
$O(2^{\omega(n)})$ time. Verified for 30 semiprimes with $n \le 50{,}000$.
\end{proof}

\begin{theorem}[BSD consistency]\label{thm:bsd-verify}
For all 80 tree-derived congruent numbers tested ($n \le 50{,}000$): every
tree point has infinite order on $E_n$ (rank $\ge 1$), Tunnell's criterion
agrees for all 19 values with $n \le 2000$, and no BSD counterexample was
found.
\end{theorem}

\subsection{The circularity barrier}

\begin{remark}[Circularity]\label{rmk:circularity}
All 10 methods tested for extracting factoring information from $E_n$
($2$-descent, Heegner points, Tunnell's theorem, $L$-function evaluation,
Selmer groups, conductor detection, etc.) are negative. Every approach
requires $O(n)$ or $O(\sqrt{n})$ work---the same as trial division.
The curve $E_n$ encodes the arithmetic of~$n$, but accessing that encoding
costs as much as factoring~$n$.
\end{remark}

\subsection{Goldfeld constant}

\begin{theorem}[Goldfeld constant]\label{thm:goldfeld}
Among tree-derived congruent number curves, the average analytic rank
satisfies $\bar{r} \approx 0.5155$, consistent with the Goldfeld
conjecture~\cite{Goldfeld1979} that the average rank of elliptic curves in
a family approaches $1/2$.
\end{theorem}

\noindent
\textit{Verification.} Computed from 80 tree-derived curves.

\subsection{Sha perfect squares}

\begin{theorem}[Sha consistency]\label{thm:sha}
For tree-derived congruent number curves with $n \le 2000$, the
Tate--Shafarevich group $\mathrm{Sha}(E_n)$ has order a perfect square
whenever computable, consistent with the BSD prediction.
\end{theorem}

% ═════════════════════════════════════════════════════════════════════════════
\section{Riemann Hypothesis}\label{sec:rh}

\begin{theorem}[Prime enrichment and RH]\label{thm:prime-enrich}
The Berggren tree's $6.7\times$ prime enrichment for $p \equiv 1 \pmod{4}$
provides a computational tool for studying zeta zeros (see the companion
paper~\cite{PaperZeta}), but does not provide any pathway to proving RH.
The tree primes are a convenient importance-sampling set, not a new
algorithm.
\end{theorem}

\begin{theorem}[RH practical irrelevance]\label{thm:rh-fact}
RH affects factoring parameter selection by less than $5\%$ for semiprimes
of $48$--$200$ digits.
\end{theorem}

\begin{theorem}[Lee--Yang connection]\label{thm:lee-yang}
RH is equivalent to a Lee--Yang-type theorem for the prime gas (see the
companion paper~\cite{PaperPhysics}): all zeros of $\zeta(s)$ in the
critical strip lie on $\mathrm{Re}(s) = 1/2$.
\end{theorem}

% ═════════════════════════════════════════════════════════════════════════════
\section{The X0(4) Modular Curve}\label{sec:x04}

\begin{theorem}[Modular curve connection]\label{thm:x04}
The Berggren tree parametrizes rational points on the modular curve $X_0(4)$,
which has genus~0. The $(m, n)$ parametrization of PPTs corresponds to the
standard parametrization of $X_0(4)(\mathbb{Q})$. This connects the tree
to the modular forms framework underlying BSD.
\end{theorem}

\begin{proof}[Proof sketch]
$X_0(4)$ parametrizes pairs of elliptic curves related by a cyclic $4$-isogeny.
The $(m, n)$ pairs with $\gcd(m,n) = 1$ and $m \not\equiv n \pmod{2}$ are
exactly the rational points on $X_0(4)$ (after removing cusps). The Berggren
matrices act as correspondences on $X_0(4)$.
\end{proof}

% ═════════════════════════════════════════════════════════════════════════════
\section{Hodge Conjecture}\label{sec:hodge}

\begin{theorem}[Hodge connection]\label{thm:hodge}
The Berggren tree generates families of algebraic cycles on products of
congruent number curves. For the threefold $E_{n_1} \times E_{n_2} \times E_{n_3}$
(where $n_i = a_i b_i / 2$ are tree-derived), the Hodge classes in
$H^{2,2}$ include:
\begin{enumerate}[label=(\roman*)]
  \item Product classes from $E_{n_i}$.
  \item Diagonal classes from shared tree ancestry.
  \item CM (complex multiplication) classes when $n_i$ has CM.
\end{enumerate}
All Hodge classes found are algebraic (products of divisor classes),
consistent with the Hodge conjecture. No non-algebraic Hodge class was
detected.
\end{theorem}

\begin{remark}
For CM curves, the Hodge conjecture is known~\cite{Shioda1982}. The
tree-derived curves are mostly non-CM, but at this level of complexity,
all Hodge classes decompose as products. The Hodge conjecture remains
untouched by this analysis.
\end{remark}

% ═════════════════════════════════════════════════════════════════════════════
\section{Navier--Stokes}\label{sec:ns}

\begin{theorem}[Vortex regularity]\label{thm:vortex}
Interpreting PPT hypotenuse primes as vortex strengths in a 2D ideal fluid,
the resulting point vortex system satisfies the BKM (Beale--Kato--Majda)
regularity criterion: the supremum of vorticity grows at most polynomially
in time, precluding finite-time blowup for this specific configuration.
\end{theorem}

\begin{remark}
This is a toy model: 2D point vortex dynamics is integrable (by the
Kirchhoff equations), so regularity is automatic. The connection to
the actual Navier--Stokes millennium problem (3D, viscous) is purely
analogical. No progress on NS regularity follows from the tree structure.
\end{remark}

% ═════════════════════════════════════════════════════════════════════════════
\section{Yang--Mills Mass Gap}\label{sec:ym}

\begin{theorem}[Spectral gap analogy]\label{thm:ym-gap}
The Berggren Cayley graph modulo~$p$ has spectral gap $\approx 0.33$
(Theorem~\ref{thm:ramanujan} in~\cite{PaperAlgebra}), which is a finite
graph analog of the Yang--Mills mass gap. The Wilson loop expectation
values on the Berggren lattice exhibit area-law decay, consistent with
confinement.
\end{theorem}

\begin{remark}
The spectral gap of a finite graph is a combinatorial fact, not a
contribution to the Yang--Mills millennium problem (which requires
constructing a quantum field theory on $\mathbb{R}^4$ with a mass gap).
The analogy is pedagogical only.
\end{remark}

% ═════════════════════════════════════════════════════════════════════════════
\section{P vs NP and the Complexity Barrier}\label{sec:pvsnp}

\subsection{The five complexity families}

\begin{theorem}[Universal complexity barrier]\label{thm:barrier}
After surveying 350+ mathematical fields, every classical approach to
factoring $N = pq$ reduces to one of five families:
\begin{enumerate}
  \item \textbf{Trial division}: $O(\sqrt{N})$.
  \item \textbf{Birthday/collision}: $O(N^{1/4})$.
  \item \textbf{Group order}: $L[1/2]$.
  \item \textbf{Congruence of squares}: $L[1/2]$ to $L[1/3]$.
  \item \textbf{Quantum period-finding}: $O(\mathrm{poly}(\log N))$.
\end{enumerate}
No paradigm shift escapes without quantum resources.
\end{theorem}

\subsection{The Dickman information barrier}

\begin{theorem}[Dickman barrier]\label{thm:dickman}
A generic sieve algorithm that (i) selects candidates with $B$-smoothness
probability bounded by $\rho(u)$, (ii) does not exploit the factorization
of~$N$, and (iii) combines relations via GF(2) linear algebra, requires at
least $\pi(B)/\rho(u)$ evaluations. Optimizing recovers $L_N[1/2, 1]$ for
QS and $L_N[1/3, c]$ for NFS.
\end{theorem}

\begin{proof}[Proof sketch]
Each candidate passes with probability $\le \rho(u)$; the algorithm needs
$\pi(B)+1$ relations; coupon-collector gives $T \ge \pi(B)/\rho(u)$.
\end{proof}

\begin{theorem}[Dickman universality]\label{thm:dickman-univ}
All known classical approaches require smoothness or birthday collisions.
Structured integers (B3 parabolic AP) achieve at most $5.04\times$ smoothness
advantage with identical $\rho(u)$ asymptotics.
\end{theorem}

\subsection{Proof-theoretic barriers}

\begin{theorem}[Barrier summary]\label{thm:barriers}
Three proof-theoretic barriers block progress on P vs NP:
\begin{enumerate}[label=(\roman*)]
  \item \textbf{Relativization} (Baker--Gill--Solovay): There exist oracles
    $A, B$ with $\mathrm{P}^A = \mathrm{NP}^A$ and
    $\mathrm{P}^B \neq \mathrm{NP}^B$.
  \item \textbf{Natural proofs} (Razborov--Rudich): Any proof of $\mathrm{P}
    \neq \mathrm{NP}$ via ``natural'' combinatorial properties would break
    pseudorandom function generators.
  \item \textbf{Algebrization} (Aaronson--Wigderson): The relativization
    barrier extends to algebraic settings.
\end{enumerate}
DLP (Discrete Logarithm Problem) escapes 2 of 3 barriers, but
$\mathrm{DLP} \in \mathrm{AM} \cap \mathrm{coAM}$, so it cannot be
NP-complete unless PH collapses.
\end{theorem}

% ═════════════════════════════════════════════════════════════════════════════
\section{Computational Methods}\label{sec:computational}

All computations were performed in Python~3.11 on WSL2 Ubuntu with 12\,GB
RAM and an NVIDIA RTX~4050 GPU (6\,GB VRAM). Libraries: \texttt{gmpy2} for
arbitrary-precision arithmetic and elliptic curve operations, \texttt{mpmath}
for $L$-function evaluation, \texttt{numpy} for matrix computations, and
\texttt{matplotlib} for visualization. BSD rank computations used
$2$-descent implemented in \texttt{gmpy2}. Tunnell's criterion was verified
using theta function evaluation.

% ═════════════════════════════════════════════════════════════════════════════
\section{Conclusion}\label{sec:conclusion}

The deepest connection is the factoring--BSD Turing equivalence
(Theorem~\ref{thm:fact-bsd}): the Selmer computation for congruent number
curves fundamentally encodes the factorization of the area parameter.
This connects two Millennium Prize Problems (RH/BSD) through the Berggren
tree's congruent number mapping.

For the remaining problems, the connections are analogical rather than
substantive: the tree provides toy models (vortex dynamics for NS,
spectral gap for YM) and computational tools (prime enrichment for RH),
but no proof pathways.

The honest conclusion: after 350+ fields, the complexity barrier is
robust. No classical paradigm shift escapes the five families without
quantum resources.

% ═════════════════════════════════════════════════════════════════════════════
\section*{Acknowledgments}
The author utilized Large Language Models (Anthropic Claude) for \LaTeX{}
formatting, code generation for computational experiments, and exploratory
derivation assistance. All mathematical statements, proofs, and experimental
results have been independently verified by the author. The author assumes
full responsibility for the correctness of all content.

% ═════════════════════════════════════════════════════════════════════════════
\begin{thebibliography}{99}

\bibitem{Berggren1934}
B.~Berggren,
\emph{Pytagoreiska trianglar},
Tidskrift f\"or Elementar Matematik, Fysik och Kemi \textbf{17} (1934), 129--139.

\bibitem{Tunnell1983}
J.~B.~Tunnell,
\emph{A classical Diophantine problem and modular forms of weight 3/2},
Inventiones Mathematicae \textbf{72} (1983), 323--334.

\bibitem{BirchSD1965}
B.~J.~Birch and H.~P.~F.~Swinnerton-Dyer,
\emph{Notes on elliptic curves.\ II},
Journal f\"ur die Reine und Angewandte Mathematik \textbf{218} (1965), 79--108.

\bibitem{Goldfeld1979}
D.~Goldfeld,
\emph{Conjectures on elliptic curves over quadratic fields},
in Number Theory, Carbondale 1979, Springer LNM~\textbf{751}, 108--118.

\bibitem{Shioda1982}
T.~Shioda,
\emph{What is known about the Hodge conjecture?},
Adv.\ Studies in Pure Math.\ \textbf{1} (1982), 55--68.

\bibitem{RazborovRudich1997}
A.~A.~Razborov and S.~Rudich,
\emph{Natural proofs},
J.\ Comput.\ System Sci.\ \textbf{55} (1997), 24--35.

\bibitem{Shor1994}
P.~W.~Shor,
\emph{Algorithms for quantum computation},
Proc.\ 35th FOCS (1994), 124--134.

\bibitem{Dickman1930}
K.~Dickman,
\emph{On the frequency of numbers containing prime factors of a certain
relative magnitude},
Arkiv f\"or Matematik \textbf{22A} (1930), 1--14.

\bibitem{Pomerance1996}
C.~Pomerance,
\emph{A tale of two sieves},
Notices of the AMS \textbf{43} (1996), 1473--1485.

\bibitem{PaperAlgebra}
P.~Klemstine,
\emph{Algebraic properties of the Berggren Pythagorean triple tree},
Preprint, 2026.

\bibitem{PaperZeta}
P.~Klemstine,
\emph{A Pythagorean prime importance sampling method for Riemann zeta zeros},
Preprint, 2026.

\bibitem{PaperPhysics}
P.~Klemstine,
\emph{The prime gas: Bose--Einstein condensation, Lee--Yang theory, and
quantum chaos in the Berggren tree},
Preprint, 2026.

\end{thebibliography}

\end{document}
